% BOOK-FILE-ID: f-d08cd421c9f0
% subsection: 等差型 \, $a_{n+1}=a_n+d$

\subsection{\texorpdfstring{等差型  $a_{n+1}=a_n+d$}{等差型}}\label{節：等差型}
まずは最もシンプルな漸化式である,等差型の漸化式を扱おう.
漸化式の一般的な形は
\[
  a_{n+1}=a_n+d
\]
として表される.
$n$が1増えるごとに$d$が1回足される.
$n$が2増えるのであれば$d$が2回足される.
したがって,その一般項を求めるには,\,$d$が何回足されたか調べて上げればよい.
漸化式を書き下して行くと,（$a_n$をより小さい$n$の状況で考えてやると）
\begin{align*}
  a_{n+1}&=a_{n}+d\\
  &=(a_{n-1}+d)+d\\
  &=(\left\{(a_{n-2}+d)\right\}+d)+d\\
  &\qquad \vdots\\
  &=a_1+nd\\
  &\therefore \quad a_{n+1}=a_1+nd
\end{align*}
$n$は任意の正整数として考えてたから,添字を全て1だけ減らして
\[
  a_n=a_1+(n-1)d
\]
が言える.この式変形が分からない人は\sectionref{節：数列と漸化式}を参照せよ.
\begin{bookpoint}
  \textbf{POINT}\par
  第$n$項に至るまでに,公差$d$は$n-1$回足される.
\end{bookpoint}


\noindent \bookblocklink{example-arithmetic-1}{problem}{answer}{【例題】}\par\nobreak
\begin{enumerate}[label=(\arabic*),parsep=3pt,format=\bookitemlink{example-arithmetic-1}{problem}{answer}]
  \item $a_1=1,\, a_{n+1}=a_n+1$
  \item $a_1=\frac{1}{2},\, a_{n+1}=a_n+\frac{1}{4}$
  \item $a_1=-\sqrt{3},\, a_{n+1}=a_n-\sqrt{2}$
\end{enumerate}

\noindent\bookblocklink{example-arithmetic-1}{hint}{problem}{【方針】}\par\nobreak
\begin{enumerate}[label=(\arabic*),parsep=3pt,format=\bookitemlink{example-arithmetic-1}{hint}{problem}]
  \item 等差型の\bookterm{general}{一般項}
    $a_n=a_1+(n-1)d$
    を用います.
  \item 公差は$d=\frac14$です.
    等差型の\bookterm{general}{一般項}に代入します.
  \item 公差は$d=-\sqrt2$です.
    負の公差の場合も同じ公式を使います.
\end{enumerate}

\noindent\bookblocklink{example-arithmetic-1}{answer}{problem}{【解答】}\par\nobreak
\begin{enumerate}[label=(\arabic*),parsep=3pt,format=\bookitemlink{example-arithmetic-1}{answer}{problem}]
  \item $a_n=1+(n-1)=n$
  \item $a_n=\frac12+\frac{n-1}{4}
      =\frac{n+1}{4}$
  \item $a_n=-\sqrt3-(n-1)\sqrt2$
\end{enumerate}

\noindent\bookblocklink{example-arithmetic-1}{solution}{problem}{【解法】}\par\nobreak
(1)は$a_1=d=1$なので$a_n=1+(n-1)=n$となる.初項は$1$,隣り合う項の差は$(n+1)-n=1$で,元の条件に合う.
(2),(3)も初項と公差を代入する.
